\begin{abstract}
Alzheimer's disease (AD) involves many biological processes at once, and researchers increasingly expect that combinations of drugs, often drugs already approved for other diseases, will be needed to treat it. Choosing which $k$ drugs to combine from $n$ candidates is a search over $\binom{n}{k}$ possibilities. We follow a network-medicine rule for good combinations: each drug's protein targets should be close to the AD disease proteins, but the drugs should not overlap with each other. We turn this rule into a single score (an ``energy'') for every combination and find the best-scoring combination in three ways. First, an \textbf{exact search} that checks every combination on the GPU gives the true best answer. Second, a \textbf{simulated-annealing search} is written three times, as sequential C++, as a hand-written AVX2 assembly kernel, and as a CUDA kernel; all three use the same random numbers and integer arithmetic, so they make exactly the same decisions, and we compare their execution times. Third, the \textbf{Quantum Approximate Optimization Algorithm (QAOA)} is run on a quantum computer simulated on the GPU with NVIDIA cuQuantum, and we measure how close it gets to the exact answer. The results are computational hypotheses, not treatment recommendations, and no quantum speedup is claimed.
\end{abstract}
